% !TEX root = chapter-standalone.tex

\section{Product of preorders}
%\linkvideo{spring2021-tradeoffs:tradeoffs:orders:composing-posets} % Composing posets
\linkvideo{spring2021-tradeoffs:tradeoffs:orders:composing-posets:product-poset} % Product of posets
Just like the product of sets, we can construct the product of \SY{preorders}.

That is a \SY{preorder} with the underlying set being the \SYN{set product}{product} of the underlying sets.

\begin{ctdefinition}[Product of preorders]
    \label{def:poset-product}
    Given \SY{preorders} $\posAdefinition$ and $\posBdefinition$, the \maindef{product preorder}
    \begin{equation}\label{eq:productposet-1}
        \posA\Ptimes \posB=\tup{\posAset \cartprod \posBset, {{\posleqof{\posA\Ptimes \posB}}}}
    \end{equation}
    is given by:

    \constit
    \begin{enumerate}
        \item The carrier set $\posAset \cartprod \posBset$;
        \item The relation $\posleqof{\posA \Ptimes \posB}$ given by
              \begin{equation}\label{eq:productposet-2}
                  \prfdoubleperiod{
                      \tup{\posAnel{1}, \posBnel{1}}
                      \posleqof{\posA\Ptimes \posB}
                      \tup{\posAnel{2}, \posBnel{2}}
                  }{
                      (\posAnel{1} \posAleq \posAnel{2})
                      \booland
                      (\posBnel{1} \posBleq \posBnel{2})
                  }
              \end{equation}
    \end{enumerate}
\end{ctdefinition}

\begin{lemma}
    \label{lem:product-preorder}
    Let~\posA and~\posB be \SY{preorders}.
    Then the relation~$\posleqof{\posA\Ptimes \posB}$ is reflexive and transitive, so~$\posA\Ptimes \posB$ is a \SY{preorder}.
    If, moreover,~\posA and~\posB are \SY{posets}, then~$\posA\Ptimes \posB$ is a \SY{poset}.
\end{lemma}
\begin{proof}
    Each property is checked one coordinate at a time.
    \begin{enumerate}
        \item Reflexivity: let~$\tup{\posela, \poselb}\setin \posAset \cartprod \posBset$.
        By reflexivity of~$\posAleq$ and of~$\posBleq$, we have~$\posela \posAleq \posela$ and~$\poselb \posBleq \poselb$, and therefore~$\tup{\posela, \poselb} \posleqof{\posA\Ptimes \posB} \tup{\posela, \poselb}$.
        \item Transitivity: suppose~$\tup{\posAnel{1}, \posBnel{1}} \posleqof{\posA\Ptimes \posB} \tup{\posAnel{2}, \posBnel{2}}$ and~$\tup{\posAnel{2}, \posBnel{2}} \posleqof{\posA\Ptimes \posB} \tup{\posAnel{3}, \posBnel{3}}$.
        By definition, this means~$\posAnel{1} \posAleq \posAnel{2}$,~$\posAnel{2} \posAleq \posAnel{3}$,~$\posBnel{1} \posBleq \posBnel{2}$, and~$\posBnel{2} \posBleq \posBnel{3}$.
        By transitivity of~$\posAleq$ and of~$\posBleq$, we get~$\posAnel{1} \posAleq \posAnel{3}$ and~$\posBnel{1} \posBleq \posBnel{3}$, that is,~$\tup{\posAnel{1}, \posBnel{1}} \posleqof{\posA\Ptimes \posB} \tup{\posAnel{3}, \posBnel{3}}$.
        \item Antisymmetry (when~\posA and~\posB are \SY{posets}): suppose~$\tup{\posAnel{1}, \posBnel{1}} \posleqof{\posA\Ptimes \posB} \tup{\posAnel{2}, \posBnel{2}}$ and~$\tup{\posAnel{2}, \posBnel{2}} \posleqof{\posA\Ptimes \posB} \tup{\posAnel{1}, \posBnel{1}}$.
        Then~$\posAnel{1} \posAleq \posAnel{2}$ and~$\posAnel{2} \posAleq \posAnel{1}$, so~$\posAnel{1} = \posAnel{2}$ by antisymmetry of~$\posAleq$; in the same way,~$\posBnel{1} = \posBnel{2}$.
        Hence~$\tup{\posAnel{1}, \posBnel{1}} = \tup{\posAnel{2}, \posBnel{2}}$.
    \end{enumerate}
\end{proof}
When~\posA and~\posB are \SY{posets}, we also call~$\posA\Ptimes \posB$ the \emph{product poset}.

Recalling the battery choice example, we have the two \SY{posets} representing cost and weight.
Given that we want to minimize both cost and weight, by considering the cost \SY{poset} containing elements ``\poscheap'', ``\posmidrange'', and ``\posexpensive'', and the weight \SY{poset} containing elements ``\poslight'', and ``\posheavy'', we can represent the product as in~\cref{fig:productpizza}.

\begin{figure*}[h!]
    \centering
    \aligninner{\includesag{70_hasse_pizza_product}}
    \caption{Product \SY{poset} of cost and weight for battery choices.}
    \label{fig:productpizza}
\end{figure*}

\begin{example}
    Consider now two \SY{posets} and their product, given in~\cref{fig:composing_posets_1}.
    \begin{figure*}[h!]
        \aligninner{
            \includesag{40_exposet_1_1}
        }
        \caption{Product of two \SY{posets}.}
        \label{fig:composing_posets_1}
    \end{figure*}
\end{example}
\vfill
\clearpage

\devel{
\subsection{Measuring the product}

\todotext{This material is DP material}


The following lemma gives expressions for the width and height of the product of two \SY{posets}.

\begin{widepar}
    \begin{lemma}
        \label{lem:measuring-product}
        If~$\posA,\posB$ are non-empty finite \SY{posets}, then we know the height of their product:
        \begin{equation}\label{eq:bound-height}
            \posetheight(\posA\Ptimes \posB) = \posetheight(\posA)+\posetheight(\posB)-1
        \end{equation}
        We can derive this bound for the width of the product:
        \begin{equation}\label{eq:bound-width}
            \posetwidth(\posA) \cdot \posetwidth(\posB)
            \leq \posetwidth(\posA\Ptimes \posB)
            \leq \min \makeset{
                \cardof \posAset \cdot \posetwidth(\posB),\cardof \posBset \cdot \posetwidth(\posA)
            }.
        \end{equation}
        This bound is tight, meaning there exist \SY{posets} for which equality holds.
    \end{lemma}
\end{widepar}

\begin{proof}
    The bound \cref{eq:bound-width} can be found in ~\cite{bezrukovantichains}.
    As for \cref{eq:bound-height}, we have the following proof.
    First, we can construct the longest \SY{chain} in~\posA:
    \begin{equation}
        \setA=\makeset{\posAnel{1},\ldots, \posAnel{\posetheight(\posA)}}.
    \end{equation}
    Furthermore, we can construct the longest \SY{chain} in~\posB:
    \begin{equation}
        \setB=\makeset{\posBel_1,\ldots, \posBnel{\posetheight(\posB)}}.
    \end{equation}
    Out of them, we can construct the \SY{chain} \begin{equation}
        \setC=\makeset{ \tup{\posAnel{1},\posBnel{1}}, \tup{\posAnel{2},\posBnel{1}},\ldots, \tup{\posAnel{\posetheight(\posA)}, \posBnel{1}}, \tup{\posAnel{\posetheight(\posA)}, \posBnel{2}},\ldots},
    \end{equation}
    which has height~$\posetheight(\posA)+\posetheight(\posB)-1$.
    So we know a lower bound for the height:
    \begin{equation}
        \posetheight(\posA\Ptimes \posB)\geq \posetheight(\posA)+\posetheight(\posB)-1.
    \end{equation}
    Now, consider a \SY{chain}~$\makeset{\tup{\posAnel{1},\posBnel{1}},\ldots, \tup{\posAnel{n},\posBnel{n}}}$ in~$\posA\Ptimes \posB$.
    By definition of chains, this means that at least one coordinate of~$\tup{\posAnel{i},\posBnel{i}}$ must be dominated by the relevant coordinate in~$\tup{\posAnel{i+1},\posBnel{i+1}}$.
    The first coordinate can only increase by $\posetheight(\posA)-1$ times, and the second one by $\posetheight(\posB)-1$ times.
    Summing up, the total number of elements in the \SY{chain} is \emph{at most} $\posetheight(\posA)+\posetheight(\posB)-1$:
    \begin{equation}
        \posetheight(\posA\Ptimes \posB)\leq \posetheight(\posA)+\posetheight(\posB)-1.
    \end{equation}
    Because upper and lower bounds are the same, we have an exact expression for the height.
    Note that this result relies on~\posA and~\posB being non-empty; if either \SY{poset} is empty, then $\posA\Ptimes \posB$ is also empty, and thus $\posetheight(\posA\Ptimes \posB)=0$.
\end{proof}
}
